\scanpage{61}
\diagram{\begin{tikzcd}[column sep=large]
(18)\quad\text{Cartes (isotop) tri-angulaires pondérées orientées}\arrow[r,"\approx"]&(\Pi_{03}\text{-Ens})\\
X\arrow[r,mapsto]&R^+(X)
\end{tikzcd}}

\margin{\uline{NB} Il n'y a pas d'autre condition sur l'ens $R$ \unsure{ayant} $\Pi_{03}$ comme groupe d'opérateurs que d'être fini\ldots}

La notion de ``carte duale'' se remplace ici par la notion de transformation d'une carte par une ``ré-indexation'' des sommets, via un
\[
g\in\mathfrak S_3\overset{\mathrm{déf}}=\mathfrak S_{\{0,1,\infty\}}.
\]
Je n'insiste pas sur ces développements, essentiellement triviaux, sauf pour noter la formule
\[
(18)\qquad R^+(X^g)=R^{\operatorname{sg}g}(X)
\qquad\begin{array}{l}g\in\mathfrak S_3\\\operatorname{sg}g\text{ la signature}\\\text{de }g\end{array}
\]

\textcircled{3} Par les considérations précédentes, et notamment par les isom (15)
\[
\begin{split}
G_2&\simeq\Pi_{03}^{\tau}/\langle\rho_1^2\rangle\\
G_2^+&\simeq\Pi_{03}/\langle\rho_1^2\rangle,
\end{split}
\]
la cat. \inserted{isotopique} des cartes finies (avec arêtes pliées) apparaît comme équivalente à une sous-catégorie pleine de la catégorie \inserted{isotop} des cartes finies tri-angulaires pondérées, et idem pour la version orientée. D'autre part, on a
\[
(19)\qquad\Pi_{03}\simeq\Pi_1(U_{03},\unsure{P_+^*|P_*^+})
\]
où \unsure{$U_{03}$:}
\[
(20)\qquad U_{03}=\mathbf P^1(\mathbf C)\setminus\{0,1,\infty\}
\]
\deleted{$P_*$ étant} $P_*$ étant un pt \deleted{quel} de l'hémisphère supérieur -- il est commode de prendre
\[
P_*=-\overline\zeta,\qquad\zeta=\exp 2i\pi/3.
\]

\mspage{60}
